\documentclass[12pt]{book}
\usepackage[T1]{fontenc}
\usepackage[brazilian]{babel}
\usepackage{csquotes}
\usepackage{makeidx}
\usepackage{glossaries}
\usepackage[
  perfil=livro,
  impressao=frente-e-verso,
  citacoes=autor-data
]{abntex3}

\makeindex
\makeglossaries
\newglossaryentry{edicao}{
  name={edição},
  description={conjunto de exemplares produzidos a partir de uma matriz}
}
\addbibresource{referencias-exemplo.bib}

\ABNTEXsetup{
  titulo={Introdução à documentação digital},
  subtitulo={conceitos, métodos e práticas},
  volume={1},
  local={Recife},
  data={2026}
}

\ABNTEXlivroSetup{
  tipo=livro,
  meio=impresso,
  editora={Editora Exemplo},
  isbn={978-65-00000-00-0},
  edicao={2},
  ano-direito-autoral={2026},
  detentor-direito-autoral={Editora Exemplo},
  direito-reproducao={Esta publicação demonstrativa pode ser reproduzida
    com indicação da fonte.},
  outros-suportes={Disponível também em formato digital},
  creditos={Projeto gráfico e revisão: Equipe Editorial Exemplo},
  dados-catalogacao={Dados demonstrativos. Substitua este conteúdo por ficha
    preparada por profissional habilitado, com nome e registro no CRB.},
  dados-editora={Rua do Livro, 10, Recife, PE; editora.example.org},
  resumo-contracapa={Uma apresentação prática dos princípios de documentação
    e preservação digital.},
  numero-paginas=120,
  capa-dura=sim,
  comporta-lombada=sim,
  possui-referencias=sim,
  possui-indice=sim
}

\ABNTEXlivroautor{Ana Silva}{autora}
\ABNTEXlivroautor{Bruno Lima}{organizador}
\ABNTEXlivrocolaborador{Carla Souza}{ilustradora}

\begin{document}

\ABNTEXlivroexterna
\ABNTEXlivrocapa
\ABNTEXlombada[Volume 1]
\ABNTEXlivrofolhasdeguarda
\ABNTEXlivroorelha
  {Ana Silva pesquisa preservação digital e normalização documental.}
  {A obra apresenta conceitos fundamentais e um fluxo demonstrativo.}

\ABNTEXlivropretextual
\ABNTEXlivrofalsafolhaderosto[Série Estudos Documentais, volume 1]
\ABNTEXlivrofolhaderosto
\ABNTEXlivrocreditos
\ABNTEXerrata
  {SILVA, Ana; LIMA, Bruno. Introdução à documentação digital. 2. ed.
   Recife: Editora Exemplo, 2026.}
  {Não há correções nesta edição demonstrativa.}
\ABNTEXdedicatoria{Às pessoas que preservam a memória documental.}
\ABNTEXagradecimentos{À equipe editorial e às pessoas revisoras.}
\ABNTEXepigrafe[Autoria demonstrativa]{Documentar é permitir reencontro.}
\ABNTEXlistadeilustracoes
\ABNTEXlistadetabelas
\ABNTEXsigla{PDF}{Portable Document Format}
\ABNTEXlistadesiglas
\ABNTEXsumario

\ABNTEXlivrotextual
\ABNTEXlivroprefacio{Este prefácio apresenta a segunda \gls{edicao}.}
\ABNTEXlivroparte{Fundamentos}
\ABNTEXlivrocapitulo{Documentação e memória}
O registro estruturado apoia a recuperação e a preservação da informação
\parencite{associacao2025}. O termo documentação\index{documentação} abrange
processos técnicos e decisões editoriais.

\begin{figure}[htbp]
  \centering
  \rule{0.55\textwidth}{3cm}
  \caption{Fluxo documental demonstrativo}
  \ABNTEXfonte{Elaborada pelas autorias (2026).}
\end{figure}

\begin{table}[htbp]
  \centering
  \caption{Componentes de uma publicação demonstrativa}
  \begin{tabular}{ll}
    \hline
    Componente & Finalidade \\
    \hline
    Metadados & Identificação editorial \\
    Estrutura & Organização do conteúdo \\
    \hline
  \end{tabular}
  \ABNTEXfonte{Elaborada pelas autorias (2026).}
\end{table}

\ABNTEXlivrocapitulo{Preservação digital}
Formatos abertos, cópias independentes e verificações de integridade reduzem
riscos de perda da informação.

\ABNTEXlivropostextual
\ABNTEXlivroposfacio{O posfácio registra possibilidades de continuidade.}
\ABNTEXbibliografia
\ABNTEXglossario
\ABNTEXappendix{Lista de verificação editorial}
Material elaborado pelas autorias para apoiar a preparação da obra.
\ABNTEXannex{Documento externo demonstrativo}
Material não elaborado pelas autorias.
\ABNTEXindice[Índice de assuntos]
\ABNTEXlivrocolofao{Composição em LaTeX com o abnTeX3 Community.\\
Editora Exemplo, Recife, 2026.}

\end{document}
